% Contenido científico íntegro normalizado del frontmatter y cortes
% Residencia material del ensamblaje sucesor de 102 capítulos.
% parte_i.tex:1--419,529--659 + parte_iii.tex:1--908.
% Se excluyen sólo la portadilla autónoma y sus órdenes de página.
\part{Correspondencias radiales genealógicas sobre el estado enriquecido: dominio, codominio y conservación de la historia helicoidal}
Esta publicación desarrolla una teoría genealógica de correspondencias
espirales desde la composición \(APP\to\TRIT\to\TPK\). El objeto primario
no es una curva presupuesta en \(\R^2\), sino una historia compatible del
estado enriquecido, junto con sus hojas aritméticas, su orientación, su
acarreo, su frontera, su ruta y su memoria. La curva aparece al final como
publicación arquimediana. En consecuencia, dos figuras iguales en
\(\R^2\) pueden proceder de historias distintas, y una circunferencia
visible puede poseer memoria radial intrínsecamente nula o ser la sombra de
una hélice cuya memoria fue eliminada por la proyección.

La obra conserva y desarrolla los resultados exactos del corpus HMT sobre
la curvatura mixta de APP, los tres regímenes cuadráticos de TRIT, la
holonomía nonádica \(\Gamma_9\), el retorno de fase con avance de memoria,
la hélice nonádica, la extensión dodecafásica no escindida y el tornillo
espectral de Pell. Sobre esos resultados construye una formalización nueva:
la envolvente afín universal, el lector radial enriquecido por prefijos, su
naturalidad bajo truncamiento, su equivariancia con \(\Gamma_9\) y su
publicación potencia--logarítmica.

La construcción no altera la teoría de masas ni reabre el electrón central.
El único punto fijo \((5,5)\) del atlas material se utiliza únicamente como
anclaje posterior de una interfaz; nunca como selector retrospectivo de la
geometría radial.

\medskip
\begin{statusbox}[title={Estratificación probatoria por dominio y codominio}]
La exposición distingue teoremas internos previamente establecidos, teoremas
demostrados en esta obra, certificaciones finitas y realizaciones físicas. Un
resultado algebraico interno se afirma con plenitud en su dominio. Toda
realización física posterior conserva sus premisas y no interviene en la
selección del generador HMT.
\end{statusbox}

\part{Construcción y resultados principales del lector radial enriquecido}
Se construye una teoría de correspondencias espirales cuyo antecedente es
el estado enriquecido de la Holografía Modular Triádica y cuya salida es una
familia de curvas en \(\R^2\) o \(\R^3\). APP proporciona un soporte
bidimensional con dos evaluaciones no simultáneamente compatibles,
\(\Sigma\) y \(\Pi\); su composición mixta posee curvatura discreta
orientada \(\pm1\). TRIT tipa los regímenes elíptico, parabólico e
hiperbólico sin confundirlos con el carácter radial ni con la curvatura de
Frenet. TPK selecciona, transporta, memoriza, actualiza y retorna el estado
completo. La holonomía \(\Gamma_9\) devuelve la fase, pero incrementa la
profundidad de memoria y transporta la frontera.

El lector inicialmente sugerido hacia
\((p,a,\Lambda,C,m,r,\theta)\) se corrige de forma decisiva. Las
coordenadas \((r,\theta)\) pertenecen a la publicación arquimediana, y el
par final \((\Lambda,C)\) no conserva sus prefijos. Se define por ello un
lector interno
\[
 \mathcal P^{\enr}_{\rad,N}:\Hist_N^{\enr}(\TPK)
 \longrightarrow \mathcal D^{\enr}_{\rad,N}
\]
que registra, en cada profundidad, los cociclos de modo y orientación, el
factor afín local, la holonomía prefija y el ledger completo. Se demuestra
su naturalidad bajo truncamiento, su extensión al límite inverso, su
equivarianza con \(\Gamma_9\), su conservación de acarreo, frontera y
memoria, y su compatibilidad con la inversión de rutas en el sector
reversible.

Después de construir el continuo se aplica un carácter radial y se introduce
la coordenada potencia--logarítmica
\[
 L_p(x)=\frac{x^p-1}{p},\qquad L_0(x)=\log x.
\]
Su ley de composición hace visible la articulación suma--producto de APP.
La traslación uniforme en \(L_p\) publica, para
\(p=2,1,0,-1,-2\), las espirales de Fermat, Arquímedes, logarítmica,
hiperbólica y lituus. La holonomía afín general produce una clasificación
más amplia mediante \((p,[\Lambda,C])\). El tornillo espectral
\(V^{-1}\mathscr SV=(2+\sqrt3)\ii\mathscr S\) aparece como caso exacto:
su proyección transversal es una espiral logarítmica y su elevación conserva
la memoria helicoidal.

La publicación concluye con la geometría diferencial de las curvas
publicadas, una realización electromagnética restringida de la hélice de
fase--posición, la interfaz con el electrón central y un sistema de
falsadores que impide confundir apariencia, genealogía y realización
física.

\bigskip
\noindent\textbf{Palabras clave:} APP; TRIT; Topological Phase Kernel;
holonomía nonádica; cociclo afín; espirales; suspensión helicoidal; memoria;
doble proyección; continuo; Mecánica Dimensional.

\part{Resultados rectores de la publicación radial APP--TRIT--TPK}
\begin{enumerate}[label=\textbf{R\arabic*.},leftmargin=12mm]
\item La geometría visible es una publicación no fiel del estado enriquecido:
igualdad en \(\R^2\) no implica igualdad de genealogía.
\item La curvatura mixta de APP distingue el orden suma--producto antes de
toda curva real.
\item \(\tau\), \(p\) y la curvatura de Frenet son invariantes de niveles
distintos y no pueden identificarse.
\item La holonomía nonádica devuelve la fase, avanza la memoria y transporta
la frontera; por tanto la proyección angular periódica admite un levantamiento
helicoidal nonádico que conserva la profundidad del retorno.
\item La información radial completa exige conservar todos los factores y
holonomías prefijas; el par final \((\Lambda,C)\) es insuficiente.
\item El lector radial enriquecido es natural bajo truncamiento y
equivariante bajo \(\Gamma_9\).
\item La involución bidireccional produce rutas conjugadas de ascenso y
descenso en el sector reversible.
\item La familia \(L_p\) publica las cinco clases potencia--logarítmicas como
salidas de una misma ley, sin elegirlas de antemano.
\item El tornillo espectral de Pell constituye una realización exacta de
rotación y autoescala cuya proyección es logarítmica.
\item Una circunferencia puede ser un estado radialmente inmóvil o la sombra
de una hélice con ledger no nulo.
\item El electrón \((5,5)\) permanece como centro material demostrado; la
firma espiral se añade como lector, no como sustituto del atlas
\(81/56/13/324\).
\item Las realizaciones electromagnéticas y materiales son posteriores a la
construcción HMT y conservan su tipado físico propio.
\end{enumerate}

\part{Notación tipada y niveles causales de la publicación radial}
\begin{longtable}{>{\raggedright\arraybackslash}p{.23\textwidth}p{.67\textwidth}}
\toprule
Símbolo & Significado\\
\midrule
\endfirsthead
\toprule
Símbolo & Significado\\
\midrule
\endhead
\(X_9\) & Toro nonádico \((\Z/9\Z)^2\), soporte bidimensional de APP.\\
\(\Sigma,\Pi\) & Evaluaciones aditiva y multiplicativa del mismo soporte.\\
\(\tau\in\{-1,0,+1\}\) & Tipo local TRIT de transporte: hiperbólico, parabólico o elíptico según la convención cuadrática declarada.\\
\(U_t\) & Transición tipada \(\mathrm{Upd}_t\circ\mathrm{Tra}_t\circ\mathrm{Sel}_t\).\\
\(\Gamma_9\) & Holonomía de nueve transiciones sobre el estado enriquecido.\\
\(p\) & Cociclo o grado radial después de definir su extracción interna; no es \(\tau\).\\
\(a\) & Carácter orientado local \(a=\tau_+-\tau_-\); se distingue del grado radial \(p=\tau_++\tau_-\).\\
\((\Lambda,C)\) & Endomorfismo afín radial: parte lineal y traslación.\\
\(H_j\) & Holonomía afín prefija hasta la profundidad \(j\).\\
\(L_p,\exp_p\) & Coordenada potencia--logarítmica y su inversa.\\
\(\kappa_{\mathrm{F}},\mathcal T_{\mathrm{F}}\) & Curvatura y torsión de Frenet de la curva publicada.\\
\(\mathcal P^{\enr}_{\rad}\) & Lector genealógico radial interno.\\
\(\ev_{\arq}^{\mathrm{esp}}\) & Publicación arquimediana de la firma espiral.\\
\bottomrule
\end{longtable}

\begin{resultbox}[title={Composición causal desde APP--TRIT--TPK hasta la evaluación arquimediana radial}]
\[
\APP\longrightarrow\TRIT\longrightarrow\TPK
\longrightarrow X_{\TPK}^{\enr}
\longrightarrow\mathcal C_{\mathrm{cont}}^{\mathrm{disc}}
\longrightarrow\mathcal P_{
  \rad}^{\enr}
\longrightarrow\ev_{\arq}^{\mathrm{esp}}
\longrightarrow\bigl(\R^2\ \text{o}\ \R^3\bigr).
\]
La curva es una salida de la cadena generativa; no interviene en la selección
del estado ni de sus coeficientes.
\end{resultbox}
\part{Fundamento genealógico de la geometría radial}

\chapter{La geometría visible como publicación del estado enriquecido}

\section{La inversión del orden explicativo}

La geometría convencional comienza por un conjunto continuo y define en él
curvas, parámetros y campos. La construcción desarrollada aquí invierte ese
orden. Primero se produce una historia compatible en un soporte finito,
después se conservan sus datos de transporte y sólo al final se publica una
coordenada real. El esquema causal es
\begin{equation}
 \APP\longrightarrow\TRIT\longrightarrow\TPK
 \longrightarrow X_{\TPK}^{\enr}
 \longrightarrow\mathcal C_{\mathrm{cont}}^{\mathrm{disc}}
 \longrightarrow\ev_{\arq}
 \longrightarrow\R^d.
 \label{espiral:eq:cadena-genealogica-espiral}
\end{equation}
No se trata de ornamentar una curva clásica con etiquetas discretas. La
historia es el objeto; la curva es una de sus lecturas.

\begin{definition}[Generación geométrica con trazabilidad]
Una construcción geométrica es \emph{constructivo--generativa} cuando:
\begin{enumerate}
\item produce su objeto desde operaciones declaradas sobre un dominio
anterior a la geometría real;
\item conserva el orden de las operaciones y las coordenadas necesarias para
reconstruirlo;
\item obtiene la apariencia geométrica mediante un mapa posterior;
\item distingue la igualdad de apariencias de la equivalencia de historias.
\end{enumerate}
\end{definition}

Este criterio es más fuerte que una parametrización. Una parametrización
\(t\mapsto\gamma(t)\) presupone la curva y cambia su modo de recorrido. Una
genealogía \(x=(x_n)_{n\ge0}\), en cambio, especifica los truncamientos
finitos, las transiciones admisibles, las hojas aritméticas, el acarreo y la
memoria que hacen posible la salida \(\gamma\).

\section{Estado, historia y publicación}

Sea \(X_{\TPK}^{\enr}\) el espacio de estados enriquecidos. Una historia
finita de longitud \(N\) es
\[
 \gamma_N=(x_0\xrightarrow{e_1}x_1\xrightarrow{e_2}\cdots
 \xrightarrow{e_N}x_N),
\]
donde cada flecha es compatible con la transición TPK tipada. Asociamos a
cada prefijo su ledger
\[
 \Led_j(\gamma_N)=
 (R_j,Q_j,\kappa_j,\sigma_j,o_j,g_j,C_j,B_j,
  \lambda_j,m_j,s_j),
\]
que registra, según la residencia concreta, residuo, cociente, acarreo,
hoja, orientación, fase, cilindro, frontera, ruta, memoria y supervivencia.
El contenido exacto puede enriquecerse, pero no puede reducirse si la
coordenada eliminada interviene en una continuación.

La publicación real es una composición
\[
 \Hist_N^{\enr}(\TPK)
 \xrightarrow{\ \mathcal P_N^{\enr}\ }
 \mathcal D_N^{\enr}
 \xrightarrow{\ \ev_N\ }
 Y,
\]
con \(Y=\R,\R^2\) o \(\R^3\). El primer mapa conserva genealogía; el
segundo publica. Si \(\ev_N(\gamma)=\ev_N(\eta)\), sólo se concluye que
ambas historias pertenecen a la misma fibra de evaluación.

\begin{proposition}[No fidelidad de la geometría visible]
Si el mapa de evaluación elimina al menos una coordenada del ledger que
admite dos valores compatibles, entonces existen historias distintas con la
misma salida geométrica.
\end{proposition}
\begin{proof}
Sea \(c\) una coordenada eliminada y sean \(\gamma,\eta\) dos historias
compatibles que coinciden en todas las coordenadas conservadas y difieren en
\(c\). Por definición de la evaluación, \(\ev(\gamma)=\ev(\eta)\), mientras
\(\gamma\neq\eta\) porque sus ledgers difieren. Por tanto \(\ev\) no es
inyectiva.
\end{proof}

La proposición no es una insuficiencia de la geometría real. Expresa su
función: contraer una estructura más rica a una coordenada utilizable. La
tesis central consiste precisamente en no confundir esa contracción con el
objeto contraído.

\section{Circunferencia intrínseca y circunferencia proyectada}

Considérese la aplicación
\[
 c_0(t)=(R\cos t,R\sin t),
\]
y la hélice
\[
 h(t)=(R\cos t,R\sin t,bt).
\]
La proyección \(\pi_{xy}(x,y,z)=(x,y)\) satisface
\(\pi_{xy}\circ h=c_0\). La misma circunferencia puede representar:
\begin{enumerate}
\item un estado con incremento radial y memoria nulos;
\item la sombra de un estado con incremento radial nulo y memoria axial
\(b\neq0\).
\end{enumerate}
La forma visible no decide entre ambos casos. Se necesita el ledger.

\begin{figure}[htbp]
\centering
\begin{subfigure}[t]{.46\textwidth}
\centering
\begin{tikzpicture}[scale=.82]
  \draw[->,HMTGray] (-2.2,0)--(2.2,0);
  \draw[->,HMTGray] (0,-2.2)--(0,2.2);
  \draw[HMTBlue,very thick] (0,0) circle (1.6);
  \fill[HMTNavy] (1.6,0) circle (1.7pt);
  \node[below,HMTGray,font=\small\sffamily] at (0,-2.25) {memoria visible nula};
\end{tikzpicture}
\caption{Circunferencia como objeto plano.}
\end{subfigure}\hfill
\begin{subfigure}[t]{.46\textwidth}
\centering
\begin{tikzpicture}[x={(.9cm,0cm)},y={(.32cm,.14cm)},z={(0cm,.72cm)},scale=.74]
  \draw[->,HMTGray] (0,0,0)--(0,0,6.8);
  \draw[HMTBlue,very thick,domain=0:720,samples=160,smooth,variable=\t]
    plot ({1.45*cos(\t)},{1.45*sin(\t)},{\t/110});
  \draw[HMTRed,dashed,thick] (0,0,0) ellipse (1.45 and .47);
  \node[right,HMTGray,font=\small\sffamily] at (0,0,6.8) {memoria};
\end{tikzpicture}
\caption{La misma sombra, con memoria levantada.}
\end{subfigure}
\caption{La geometría visible no reconstruye por sí sola la genealogía.}
\label{espiral:fig:circulo-helice}
\end{figure}

\begin{resultbox}[title={Defecto de fidelidad de la evaluación arquimediana sobre clases genealógicas}]
Una circunferencia puede poseer memoria intrínsecamente nula y puede también
ser la publicación de una historia helicoidal. La igualdad de la curva no
autoriza a identificar los estados que la producen.
\end{resultbox}

\chapter{APP: soporte bidimensional, hojas aritméticas y curvatura mixta}

\section{La célula nonádica y el nacimiento bidimensional}

La célula visible
\[
 E_9=\{1,2,3,4,5,6,7,8,9\}
\]
porta simultáneamente una fase y una carta de representantes positivos. Para
\(n\ge1\),
\[
 \rho_9(n)=1+((n-1)\bmod9),\qquad
 q_9(n)=\left\lfloor\frac{n-1}{9}\right\rfloor,
\]
de modo que
\[
 n=9q_9(n)+\rho_9(n).
\]
El producto cartesiano de dos células construye
\[
 X_9=(\Z/9\Z)^2.
\]
No es una recta repetida, sino un soporte genuinamente bidimensional. La
completación \(9\times9\) puede leerse como un núcleo \(8\times8\) y un
gnomon de \(17\) celdas:
\[
81=64+17.
\]

Sobre cada \((i,j)\in X_9\) actúan dos evaluaciones del mismo soporte:
\[
 \Sigma(i,j)=\operatorname{dr}(i+j),\qquad
 \Pi(i,j)=\operatorname{dr}(ij),
\]
donde \(\operatorname{dr}\) publica el representante nonádico y el
levantamiento conserva el cociente. APP no suma dos espacios; conserva un
solo espacio con dos lecturas.

\section{Residuo, cociente y acarreo}

Las identidades levantadas
\[
 i+j=R^+_{ij}+9q^+_{ij},\qquad
 ij=R^\times_{ij}+9q^\times_{ij}
\]
distinguen la salida visible \(R\) de la elevación \(q\). Si se conserva
únicamente el residuo, operaciones distintas pueden colapsar en la misma
celda. La pareja \((R,q)\) despliega la historia. Bajo composición, el
acarreo obedece una ley de cociclo; por ello no es una anotación auxiliar,
sino el dato que permite recomponer una operación larga desde sus pasos.

\begin{example}[Igual residuo, distinta historia]
En la carta nonádica,
\[
 8+8=7+9\cdot1,\qquad 8\cdot2=7+9\cdot1.
\]
Las dos operaciones publican el mismo residuo y el mismo cociente local,
pero difieren en hoja. La etiqueta \(\Sigma/\Pi\) continúa siendo necesaria.
En rutas más largas, el orden de esas hojas modifica el acarreo compuesto y
la frontera alcanzada.
\end{example}

\section{Incompatibilidad operatoria de las hojas aditiva y multiplicativa}

Para \(d\ge2\), sea \(\mathscr A_d=E_9^d\). Las esperanzas condicionales
\(P_{\Sigma,d}\) y \(P_{\Pi,d}\) sobre las dos evaluaciones satisfacen
\[
 \|[P_{\Sigma,2},P_{\Pi,2}]\|=\frac{\sqrt2}{3},
 \qquad
 \|[P_{\Sigma,3},P_{\Pi,3}]\|=\frac{\sqrt3}{4}.
\]
En dimensión tres, el modo trítico de fase operativa posee valor singular
\(s=1/2\), y
\[
 s^2=\frac14,\qquad 1-s^2=\frac34=\frac{90}{120},
 \qquad s\sqrt{1-s^2}=\frac{\sqrt3}{4}.
\]
Así, la separación entre las hojas no se añade después: está cuantificada en
el propio soporte APP.

\section{Curvatura discreta suma--producto}

En residuos, escribamos
\[
 S(x,y)=x+y,\qquad P(x,y)=xy\pmod9.
\]
Para un potencial \(T\), se definen las diferencias progresivas
\[
 \Delta_xT(x,y)=T(x+1,y)-T(x,y),\qquad
 \Delta_yT(x,y)=T(x,y+1)-T(x,y).
\]
Con \(A_x^T=\Delta_xT\) y \(A_y^T=\Delta_yT\), la curvatura de una mezcla
ordenada es
\[
 F(T_x,T_y)=\Delta_xA_y^{T_y}-\Delta_yA_x^{T_x}.
\]
Como
\[
 \Delta_xS=\Delta_yS=1,\qquad
 \Delta_xP=y,\qquad \Delta_yP=x,
\]
se obtiene exactamente
\begin{equation}
 F(S,S)=F(P,P)=0,\qquad
 F(S,P)=+1,\qquad F(P,S)=-1.
 \label{espiral:eq:curvatura-mixta-app}
\end{equation}

\begin{theorem}[Curvatura mixta orientada de APP]
Las composiciones puras de las hojas de APP son planas en la convención de
plaquetas precedente; las composiciones mixtas poseen curvatura constante de
signo opuesto según el orden.
\end{theorem}
\begin{proof}
Sustituyendo las cuatro diferencias en la definición,
\[
 \Delta_x\Delta_yS-\Delta_y\Delta_xS=0,
 \quad
 \Delta_x\Delta_yP-\Delta_y\Delta_xP=0
\]
para las dos composiciones puras, mientras
\[
 \Delta_x\Delta_yP-\Delta_y\Delta_xS=1,
 \qquad
 \Delta_x\Delta_yS-\Delta_y\Delta_xP=-1.
\]
La identidad se verifica en cada una de las \(81\) plaquetas.
\end{proof}

Para un rectángulo coordenado de lados \(a,b\), la cancelación de las
aristas interiores da la forma finita de Stokes:
\[
 W_{\partial R_{a,b}}(S,P)=ab,\qquad
 W_{\partial R_{a,b}}(P,S)=-ab\pmod9.
\]

\begin{figure}[htbp]
\centering
\begin{tikzpicture}[scale=.72]
\foreach \x in {0,...,8}{
  \draw[HMTLine] (\x,0)--(\x,8);
  \draw[HMTLine] (0,\x)--(8,\x);
}
\draw[HMTNavy,thick] (0,0) rectangle (8,8);
\draw[->,HMTBlue,very thick] (2,2)--(5,2) node[midway,below] {$S$};
\draw[->,HMTRed,very thick] (5,2)--(5,5) node[midway,right] {$P$};
\draw[->,HMTBlue!70!black,very thick] (5,5)--(2,5);
\draw[->,HMTRed!80!black,very thick] (2,5)--(2,2);
\node[fill=white,inner sep=2pt] at (3.5,3.5) {$F(S,P)=+1$};
\node[below,HMTGray,font=\small\sffamily] at (4,-.75) {soporte APP \(9\times9\)};
\end{tikzpicture}
\caption{La curvatura nace del orden de las hojas sobre un mismo soporte.}
\label{espiral:fig:curvatura-app}
\end{figure}

\section{Antecedentes APP de la clasificación radial}

APP aporta tres ingredientes distintos:
\begin{enumerate}
\item un dominio bidimensional anterior a \(\R^2\);
\item dos operaciones cuya mezcla ordenada distingue orientación;
\item un ledger de residuo, cociente, hoja y acarreo capaz de recordar el
camino que produjo una misma salida.
\end{enumerate}
La curvatura \(\pm1\) es el antecedente local de la orientación radial, pero
no debe identificarse sin prueba con el grado global \(p\). El paso de uno a
otro requiere la acumulación TPK y el cociclo declarado en la
\cref{espiral:chap:cociclos-radiales}.

\chapter{TRIT: orientación local y 3 regímenes geométricos}

\section{Estado ternario levantado}

El observable visible es
\[
 \TRIT=\{-1,0,+1\}.
\]
En un paso local con un acarreo, la descomposición balanceada se escribe
\[
 \widehat\tau=(r,c),\qquad a+b=r+3c,
 \qquad r\in\{-1,0,+1\},\ c\in\Z.
\]
El cero visible posee al menos tres levantamientos elementales:
\[
 0^+=(0,+1),\qquad0^0=(0,0),\qquad0^-=(0,-1).
\]
La igualdad visible no fija la continuación. Ésta depende también de hoja,
orientación, altura, acarreo, frontera, ruta y fase.

\section{Álgebras cuadráticas}

Para \(\tau\in\TRIT\), sea
\[
 \mathbb A_\tau=\R[J]/(J^2+\tau I).
\]
Se obtienen tres regímenes:
\[
 J_{+1}^2=-I,\qquad J_0^2=0,\qquad J_{-1}^2=+I,
\]
y una sola identidad exponencial
\[
 \exp(tJ_\tau)=C_\tau(t)I+S_\tau(t)J_\tau,
\]
donde
\[
\begin{array}{c|c|c|c}
\tau & J_\tau^2 & C_\tau(t)&S_\tau(t)\\ \hline
+1&-I&\cos t&\sin t\\
0&0&1&t\\
-1&I&\cosh t&\sinh t
\end{array}
\]

\begin{figure}[htbp]
\centering
\begin{tikzpicture}[>=Latex]
\begin{scope}[xshift=-5cm]
 \draw[->,HMTGray] (-1.8,0)--(1.8,0);
 \draw[->,HMTGray] (0,-1.8)--(0,1.8);
 \draw[HMTBlue,very thick] (0,0) circle (1.25);
 \draw[->,HMTRed,thick] (0,0)--({1.25*cos(40)},{1.25*sin(40)});
 \node[below,HMTNavy,font=\sffamily\bfseries] at (0,-2.05) {elíptico \(J^2=-I\)};
\end{scope}
\begin{scope}
 \draw[->,HMTGray] (-1.8,0)--(1.8,0);
 \draw[->,HMTGray] (0,-1.8)--(0,1.8);
 \draw[HMTBlue,very thick] (-1.3,-.65)--(1.3,.65);
 \draw[->,HMTRed,thick] (0,0)--(1.1,.55);
 \node[below,HMTNavy,font=\sffamily\bfseries] at (0,-2.05) {parabólico \(J^2=0\)};
\end{scope}
\begin{scope}[xshift=5cm]
 \draw[->,HMTGray] (-1.8,0)--(1.8,0);
 \draw[->,HMTGray] (0,-1.8)--(0,1.8);
 \draw[HMTBlue,very thick,domain=-1.3:1.3,samples=80] plot ({cosh(\x)},{sinh(\x)});
 \draw[HMTBlue,very thick,domain=-1.3:1.3,samples=80] plot ({-cosh(\x)},{sinh(\x)});
 \draw[->,HMTRed,thick] (0,0)--(1.25,.75);
 \node[below,HMTNavy,font=\sffamily\bfseries] at (0,-2.05) {hiperbólico \(J^2=I\)};
\end{scope}
\end{tikzpicture}
\caption{Los tres regímenes cuadráticos del TRIT. Son tipos de transporte, no nombres de curvas concretas.}
\label{espiral:fig:trit-regimenes}
\end{figure}

\section{Separación entre régimen TRIT, carácter radial y curvatura de Frenet}

La teoría requiere tres niveles que no deben aplanarse:
\begin{enumerate}
\item \(\tau\) tipa el régimen local de transporte;
\item \(p\) tipa el carácter radial acumulado de una historia;
\item \(\kappa_{\mathrm{F}}\) y \(\mathcal T_{\mathrm{F}}\) describen la curva ya
publicada en un espacio real.
\end{enumerate}
Una circunferencia posee holonomía angular elíptica, pero el presente TRIT
\(0\) es un umbral parabólico. La circunferencia y el presente no son el
mismo objeto. Del mismo modo, una espiral logarítmica publicada puede incluir
tramos cuyas transiciones locales recorren los tres tipos TRIT.

\begin{principle}[Separación de niveles]
Ninguna identidad entre \(\tau\), \(p\) y \(\kappa_{\mathrm{F}}\) se admite por
semejanza de signo, nombre o dibujo. Toda relación debe construirse mediante
un mapa con dominio, codominio y acción declarados.
\end{principle}

\section{Involución de orientación}

Para una palabra \(\boldsymbol\tau=(\tau_1,\ldots,\tau_n)\),
\[
 C_{\mathrm{rev}}(\boldsymbol\tau)
 =(-\tau_n,\ldots,-\tau_1),
 \qquad C_{\mathrm{rev}}^2=\id.
\]
La involución intercambia las orientaciones y preserva el umbral. En el
sector de rutas reversibles inducirá la inversión de la holonomía afín, pero
esa conclusión requiere que la representación local satisfaga
\(h_{e^\dagger}=h_e^{-1}\).

\chapter{Estado enriquecido, continuo único y doble proyección}

\section{Sistema inverso de historias}

Sean \(\mathsf S_n\) los estados admisibles a profundidad \(n\) y
\[
 \pi_{n+1,n}:\mathsf S_{n+1}\longrightarrow\mathsf S_n
\]
los truncamientos. El espacio de historias completas es
\[
 X_\infty=\varprojlim_n\mathsf S_n.
\]
Con la ultramétrica
\[
 d(x,y)=3^{-m(x,y)},
\]
donde \(m(x,y)\) es la primera profundidad en que difieren, los cilindros de
prefijo fijo son abiertos y cerrados. La estructura del continuo se obtiene
al construir conjuntamente, desde el mismo estado APP--TRIT--TPK, sus cinco
construcciones consustanciales, que emergen correlativamente, y pasar por terminalidad, naturalidad y
límite. Esta publicación no las disgrega ni las redefine mediante curvas.

\section{Evaluación arquimediana mediante cilindros compatibles}

Sea \(\widehat{\mathbb Q}\) la completación uniforme de \(\mathbb Q\), construida como
clases de sucesiones racionales de Cauchy. La construcción calibrada asocia
a cada estado finito una dirección entera \(k_n(s)\) y el cilindro racional
semicerrado
\[
 I_n(s)=\{q\in\mathbb Q:k_n(s)3^{-n}\le q<(k_n(s)+1)3^{-n}\}.
\]
Satisface simultáneamente:
\begin{enumerate}[label=\textup{(A\arabic*)}]
\item \(\bigcup_{s\in\mathsf S_n}I_n(s)=\mathbb Q\) para todo \(n\);
\item si \(t\in\mathsf S_{n+1}\) trunca a \(s\), entonces
      \(I_{n+1}(t)\subseteq I_n(s)\), y los hijos particionan el padre con
      la convención semicerrada de frontera;
\item existe una constante \(C>0\) y \(0<\lambda<1\) tales que
\[
 \operatorname{diam}I_n(s)\le C\lambda^n,
 \qquad 0<\lambda<1;
\]
\item toda cadena de cilindros anidados procede de una historia compatible,
      con las dos expansiones de frontera identificadas por la relación
      arquimediana.
\end{enumerate}

\begin{theorem}[Cierre arquimediano del continuo genealógico]
\label{thm:espiral-cierre-arquimediano-continuo-genealogico}
Cada historia calibrada determina un único valor
\[
 \{\operatorname{val}(x)\}
 =\bigcap_n\overline{I_n(x_n)}^{\widehat{\mathbb Q}}.
\]
La evaluación
\[
 \operatorname{val}:X_\infty^{\mathrm{cal}}\longrightarrow\widehat{\mathbb Q}
\]
es continua, Hölder y sobreyectiva. El cociente por igualdad de valor
satisface
\[
 X_\infty^{\mathrm{cal}}/{\sim_{\arq}}
 \cong\widehat{\mathbb Q}\cong\R,
\]
donde la última identificación es el reconocimiento arquimediano usual de
la completación de \(\mathbb Q\), no una entrada del generador.
\end{theorem}
\begin{proof}
La contracción de diámetros y la completitud de \(\widehat{\mathbb Q}\) reducen cada
cadena anidada a un punto. La compatibilidad de cilindros demuestra continuidad; la
cota geométrica de diámetros da la estimación de Hölder. Por (A1) y (A2),
cada clase de Cauchy racional selecciona a toda profundidad un cilindro que
la contiene;
(A4) proporciona una historia compatible y resuelve de modo único las
fronteras. De aquí resulta la sobreyectividad. Dos historias tienen la misma
imagen exactamente cuando pertenecen a la relación \(\sim_{\arq}\), por lo
que la aplicación inducida en el cociente es una biyección continua y abierta
sobre la topología final de cilindros; así es un homeomorfismo.
\end{proof}

La recta real es, por tanto, una salida demostrada de la evaluación y no un
dominio presupuesto. La historia no se reduce a su valor.

\section{Doble proyección}

Desde un mismo dominio calibrado se distinguen
\[
 \mathfrak R:X_\infty^{\mathrm{cal}}\longrightarrow\R,
 \qquad
 \mathfrak E:X_\infty^{\mathrm{cal}}\times\mathscr C_{\mathrm{exc}}
 \longrightarrow\mathbf{Exc}.
\]
La primera contrae cilindros compatibles; la segunda conserva incidencia de
frontera hacia el corredor excepcional. La teoría espiral se inserta en la
primera rama, pero conserva en su ledger los datos necesarios para comparar
la segunda. No se identifica el valor con la incidencia ni se hace que una
publicación genere retrospectivamente el continuo.

\begin{figure}[htbp]
\centering
\begin{tikzpicture}[>=Latex,node distance=12mm and 13mm,
 n/.style={draw=HMTNavy,fill=white,minimum height=9mm,text width=34mm,align=center,font=\small\sffamily}]
\node[n] (x) {$X_\infty^{\mathrm{cal}}$\\estado enriquecido común};
\node[n,below left=of x] (r) {publicación arquimediana\\valor y curva en \(\R^d\)};
\node[n,below right=of x] (e) {publicación excepcional\\incidencia y automorfismos};
\draw[->,HMTBlue,very thick] (x)--node[left,font=\small] {$\mathfrak R$} (r);
\draw[->,HMTRed,very thick] (x)--node[right,font=\small] {$\mathfrak E_c$} (e);
\draw[HMTGold,dashed,thick] (r)--node[below,font=\scriptsize\sffamily]{invariantes comunes} (e);
\end{tikzpicture}
\caption{Las dos proyecciones son correlativas y posteriores al mismo estado.}
\label{espiral:fig:doble-proyeccion-espiral}
\end{figure}

\section{5 regiones de \texorpdfstring{\(\pi\)}{pi} y multiplicidad de genealogías}

La microfibra de \(\pi\) contiene cinco regiones distintas con
multiplicidades
\[
 432+4\cdot144=1008.
\]
El mismo valor visible conserva historias regionales diferentes. Este hecho
proporciona el precedente exacto de la clasificación espiral: igualdad de
una curva publicada no obliga a igualdad de hoja, carry, frontera ni
memoria. La teoría no inventa esa multiplicidad; la transporta a un nuevo
tipo de salida geométrica.

\begin{resultbox}[title={Antecedencia causal del continuo discreto respecto del lector radial}]
Hasta aquí se han construido el dominio, las hojas, los regímenes, la
holonomía, el estado enriquecido y la publicación del continuo. Sólo ahora
es lícito introducir coordenadas reales \((r,\theta)\) y reconocer curvas
convencionales.
\end{resultbox}
\part{Teoría genealógica de correspondencias espirales}

\chapter{Envolvente afín universal de las hojas APP}

\section{Construcción del módulo radial universal y de su envolvente endomórfica}

Las operaciones \(u\mapsto u+a\) y \(u\mapsto\lambda u\) sugieren el
grupo afín de la recta. Sin embargo, APP contiene sectores multiplicativos
no invertibles y precede a la publicación real. El blanco interno apropiado
es por ello el monoide de endomorfismos afines de un módulo radial universal.

Sea \(\widetilde X_{\APP}\) el conjunto de estados APP completos, es decir,
estados que conservan ambas hojas, residuo, cociente y acarreo, y añadamos un
símbolo absorbente \(\bot\). Definimos el módulo libre
\begin{equation}
 U_{\APP}:=\Z^{(\widetilde X_{\APP}\sqcup\{\bot\})}/\langle[\bot]\rangle.
 \label{espiral:eq:modulo-radial-universal}
\end{equation}
Para cada transición enriquecida \(e\), de fuente \(x_e\) y dígitos activos
\((i_e,j_e)\), APP produce sin evaluación decimal los registros
\[
 i_e+j_e=R_e^++9q_e^+,
 \qquad
 i_ej_e=R_e^\times+9q_e^\times.
\]
Sean \(s_e^\Sigma(x_e)\) y \(s_e^\Pi(x_e)\) los estados completos obtenidos
al adjuntar, respectivamente, \((R_e^+,q_e^+)\) y
\((R_e^\times,q_e^\times)\), conservando la otra hoja en el ledger. El
levantamiento aditivo determina el vector explícito
\[
 A_e=[s_e^\Sigma(x_e)]-[x_e]\in U_{\APP},
\]
y el levantamiento multiplicativo determina la aplicación total
\(m_e:\widetilde X_{\APP}\sqcup\{\bot\}\to
\widetilde X_{\APP}\sqcup\{\bot\}\), prolongada por
\[
 m_e(x)=
 \begin{cases}
 s_e^\Pi(x_e),&x=x_e,\\
 \bot,&x\neq x_e,
 \end{cases}
 \qquad m_e(\bot)=\bot.
\]
Su linealización es
\[
 M_e([x])=[m_e(x)]\in U_{\APP}.
\]
Así, \(A_e\) y \(M_e\) no son parámetros libres: son las dos publicaciones
linealizadas de la misma transición APP completa.

Definimos
\[
 \AffEnd(U_{\APP})
 =\operatorname{End}(U_{\APP})\ltimes U_{\APP}
\]
con composición
\begin{equation}
 (\Lambda_2,C_2)\circ(\Lambda_1,C_1)
 =\bigl(\Lambda_2\Lambda_1,
 C_2+\Lambda_2C_1\bigr).
 \label{espiral:eq:producto-affend}
\end{equation}
El par actúa por \(u\mapsto\Lambda u+C\). La parte lineal puede ser no
invertible. El subgrupoide de unidades se obtiene restringiendo
\(\Lambda\in\operatorname{Aut}(U_{\APP})\).

\begin{proposition}[Propiedad universal]
Toda aplicación del conjunto de estados completos a un módulo \(V\), nula
en \(\bot\), se prolonga de manera única a un morfismo
\(\chi:U_{\APP}\to V\). Si además entrelaza los levantamientos
multiplicativos, \(\chi M_e=M_e^V\chi\), entonces transporta de manera única
los factores afines \(h_e\) y todas sus holonomías prefijas.
\end{proposition}
\begin{proof}
La primera afirmación es la propiedad universal del módulo libre cocientado
por \([\bot]=0\). La segunda sigue aplicando \(\chi\) a la composición
semidirecta: \(\chi(C_2+\Lambda_2C_1)=
\chi(C_2)+\Lambda_2^V\chi(C_1)\).
\end{proof}

\section{Representación local de una transición}

Sea \(e:x\to y\) una transición TPK. El selector trítico de fase operativa
determina si actúa la hoja aditiva, la multiplicativa o el umbral. Se define
\begin{equation}
 h_e=
 \begin{cases}
 (I,A_e),&\text{hoja }\Sigma,\\
 (M_e,0),&\text{hoja }\Pi,\\
 (I,0),&\text{umbral},
 \end{cases}
 \label{espiral:eq:factor-afin-local}
\end{equation}
donde \(A_e\) y \(M_e\) son los levantamientos APP completos construidos
arriba. No se toman sólo sus residuos: el cociente y el acarreo permanecen
en el ledger.

\begin{definition}[Representación radial interna]
Una representación radial interna es una asignación
\[
 \rho_{\rad}:\mathscr P_{\TPK}^{\enr}
 \longrightarrow\AffEnd(U_{\APP})
\]
que envía cada flecha a \(h_e\), la identidad a \((I,0)\) y la composición
de caminos al producto \cref{espiral:eq:producto-affend}.
\end{definition}

La definición no consulta una curva. Los factores proceden de la hoja y del
levantamiento que el estado ya contiene. Una realización numérica posterior
puede enviar \(U_{\APP}\) a \(\R\), pero esa evaluación no modifica la
representación interna.

\section{Holonomías prefijas}

Para
\[
 \gamma_N=e_N\cdots e_1,
\]
definimos
\[
 H_0=(I,0),\qquad
 H_j=h_{e_j}\cdots h_{e_1}=(\Lambda_j,C_j).
\]

\begin{theorem}[Holonomía afín functorial]
Para rutas componibles \(\gamma_1,\gamma_2\),
\[
 H_{\gamma_2\gamma_1}=H_{\gamma_2}H_{\gamma_1}.
\]
Si \(H_{\gamma_i}=(\Lambda_i,C_i)\), entonces
\begin{equation}
 \Lambda_{\gamma_2\gamma_1}=\Lambda_2\Lambda_1,
 \qquad
 C_{\gamma_2\gamma_1}=C_2+\Lambda_2C_1.
 \label{espiral:eq:holonomia-afin-functorial}
\end{equation}
\end{theorem}
\begin{proof}
La acción sucesiva sobre \(u\) es
\[
 u\xmapsto{H_{\gamma_1}}\Lambda_1u+C_1
 \xmapsto{H_{\gamma_2}}
 \Lambda_2\Lambda_1u+C_2+\Lambda_2C_1.
\]
La unicidad de la descomposición afín da \cref{espiral:eq:holonomia-afin-functorial}.
\end{proof}

Para un bloque de nueve pasos,
\begin{align}
 \Lambda_9&=\Lambda_{e_9}\cdots\Lambda_{e_1},\\
 C_9&=C_{e_9}+\Lambda_{e_9}C_{e_8}
 +\Lambda_{e_9}\Lambda_{e_8}C_{e_7}+\cdots
 +\Lambda_{e_9}\cdots\Lambda_{e_2}C_{e_1}.
 \label{espiral:eq:traslacion-nueve}
\end{align}
La posición de cada factor importa. Ésta es la forma afín de la memoria de
orden.

\section{Conmutador suma--producto}

En el sector invertible y sobre una realización unidimensional, sean
\[
 A_a(u)=u+a,\qquad M_\lambda(u)=\lambda u.
\]
Entonces
\begin{equation}
 A_aM_\lambda A_a^{-1}M_\lambda^{-1}
 =A_{a(1-\lambda)}.
 \label{espiral:eq:conmutador-afin}
\end{equation}
La prueba es directa:
\[
 u\mapsto\lambda^{-1}u\mapsto\lambda^{-1}u-a
 \mapsto u-\lambda a\mapsto u+a(1-\lambda).
\]
La identidad sólo se afirma cuando \(\lambda\) es una unidad. En el monoide
completo, el defecto de conmutación se conserva mediante los factores y el
ledger, no mediante un conmutador que requiera inversas inexistentes.

\begin{figure}[!htb]
\centering
\begin{tikzpicture}[>=Latex,node distance=18mm and 28mm,
 n/.style={draw=HMTNavy,fill=white,minimum width=23mm,minimum height=8mm,font=\small}]
\node[n] (u) {$u$};
\node[n,right=of u,yshift=10mm] (a) {$u+a$};
\node[n,right=of a] (ma) {$\lambda u+\lambda a$};
\node[n,right=of u,yshift=-10mm] (m) {$\lambda u$};
\node[n,right=of m] (am) {$\lambda u+a$};
\draw[->,HMTBlue,thick] (u)--node[above left] {$A_a$} (a);
\draw[->,HMTRed,thick] (a)--node[above] {$M_\lambda$} (ma);
\draw[->,HMTRed,thick] (u)--node[below left] {$M_\lambda$} (m);
\draw[->,HMTBlue,thick] (m)--node[below] {$A_a$} (am);
\draw[->,HMTGold,dashed,thick] (ma)--node[right,font=\scriptsize]
  {$A_{a(1-\lambda)}$} (am);
\node[below=12mm of u,HMTGray,font=\small\sffamily,text width=80mm,align=center]
{Los dos órdenes terminan en puntos distintos; la flecha vertical mide
exactamente el defecto de composición.};
\end{tikzpicture}
\caption{La diferencia suma--producto se manifiesta como holonomía afín.}
\label{espiral:fig:cuadrado-afin}
\end{figure}

\chapter{Cociclos radiales en el levantamiento TRIT de 2 semicoronas}
\label{espiral:chap:cociclos-radiales}

\section{Categoría levantada y caracteres elementales}

Sea \(\mathscr P_{\TPK}^{2\mathrm c,\enr}\) la categoría de caminos cuyas
flechas son ternas
\[
 \widehat e=(e,\tau_+(e),\tau_-(e)),
 \qquad \tau_\pm(e)\in\{-1,0,+1\},
\]
donde el par orientado pertenece al estado enriquecido y el funtor de olvido
envía \(\widehat e\) a la transición TPK \(e\). Definimos en cada generador
\begin{equation}
 p(\widehat e)=\tau_+(e)+\tau_-(e),
 \qquad
 a(\widehat e)=\tau_+(e)-\tau_-(e).
 \label{espiral:eq:carta-doble-trit}
\end{equation}
Para una ruta \(\widehat\gamma_N=\widehat e_N\cdots\widehat e_1\), los
cociclos acumulados son
\begin{equation}
 P_j=\sum_{k=1}^{j}p(\widehat e_k),
 \qquad
 A_j=\sum_{k=1}^{j}a(\widehat e_k),
 \qquad P_0=A_0=0.
 \label{espiral:eq:cociclos-radiales}
\end{equation}
La aditividad bajo concatenación demuestra que \((P,A)\) es un cociclo
\(\Z^2\)-valuado de la categoría levantada. Esta afirmación no identifica por
decreto el par con otro cociclo previamente nombrado: lo construye a partir
de las dos orientaciones efectivamente conservadas en cada flecha. La
subdivisión por microflechas quietas preserva ambas sumas.

\section{Carta local exhaustiva de los caracteres radial y orientado}

\begin{longtable}{cccccc}
\caption{Carta local completa de las dos orientaciones TRIT.}
\label{espiral:tab:doble-trit}\\
\toprule
\(\tau_+\)&\(\tau_-\)&\(p\)&\(a\)&\(r_3(p)\)&\(c_3(p)\)\\
\midrule
\endfirsthead
\toprule
\(\tau_+\)&\(\tau_-\)&\(p\)&\(a\)&\(r_3(p)\)&\(c_3(p)\)\\
\midrule
\endhead
+1&+1&+2&0&-1&1\\
+1&0&+1&+1&+1&0\\
+1&-1&0&+2&0&0\\
0&+1&+1&-1&+1&0\\
0&0&0&0&0&0\\
0&-1&-1&+1&-1&0\\
-1&+1&0&-2&0&0\\
-1&0&-1&-1&-1&0\\
-1&-1&-2&0&+1&-1\\
\bottomrule
\end{longtable}

Los tres estados con \(p=0\) se distinguen por
\(a=2,0,-2\). Por tanto el grado radial no basta para
reconstruir la orientación.

\section{Descomposición balanceada}

Para todo \(p\in\Z\), existe una única descomposición
\[
 p=r_3(p)+3c_3(p),
 \qquad r_3(p)\in\{-1,0,+1\},\quad c_3(p)\in\Z.
\]
En el bloque local \(p\in[-2,2]\), la última columna de
\cref{espiral:tab:doble-trit} da el acarreo. En rutas largas, \(c_3(p)\) no queda
restringido a \(\{-1,0,+1\}\). Esta precisión evita convertir un ejemplo
local en una ley global falsa.

\section{Involución}

La involución local
\[
 (\tau_+,\tau_-)\longmapsto(-\tau_-,-\tau_+)
\]
produce
\[
 p\longmapsto-p,
 \qquad
 a\longmapsto a.
\]
El carácter radial cambia de orientación, mientras la diferencia orientada
permanece. Esta separación será visible en la identidad
\(L_{-p}(x^{-1})=-L_p(x)\).

\chapter{Coordenada potencia--logarítmica y ley suma--producto}

\section{Definición y dominio}

Para \(x>0\), definimos
\begin{equation}
 L_p(x)=
 \begin{cases}
 \dfrac{x^p-1}{p},&p\neq0,\\[2mm]
 \log x,&p=0.
 \end{cases}
 \label{espiral:eq:lp-def}
\end{equation}
Su inversa es
\begin{equation}
 \exp_p(u)=
 \begin{cases}
 (1+pu)^{1/p},&p\neq0,\quad1+pu>0,\\[1mm]
 \e^u,&p=0.
 \end{cases}
 \label{espiral:eq:exp-p}
\end{equation}
La condición \(1+pu>0\) no es ornamental: fija la rama real positiva.

\section{Adición deformada}

En el dominio correspondiente, sea
\[
 u\boxplus_pv=u+v+puv.
\]

\begin{theorem}[Ley suma--producto]
Para \(x,y>0\),
\begin{equation}
 L_p(xy)=L_p(x)\boxplus_pL_p(y).
 \label{espiral:eq:lp-producto}
\end{equation}
Además, \(0\) es la identidad de \(\boxplus_p\), y
\[
 \exp_p(u\boxplus_pv)=\exp_p(u)\exp_p(v).
\]
\end{theorem}
\begin{proof}
Para \(p\neq0\),
\[
 1+pL_p(xy)=(xy)^p=x^py^p
 =(1+pL_p(x))(1+pL_p(y)).
\]
Al expandir el último producto y dividir por \(p\) se obtiene
\cref{espiral:eq:lp-producto}. Para \(p=0\) es la identidad
\(\log(xy)=\log x+\log y\). La igualdad para \(\exp_p\) sigue al aplicar
la inversa.
\end{proof}

Esta ley expresa de forma exacta la articulación APP: el producto en la
coordenada radial se publica como una suma deformada en la coordenada
\(L_p\). El término \(puv\) conserva la interacción que desaparecería en
una suma ordinaria.

\section{Inversión bidireccional}

\begin{proposition}[Identidad de inversión]
Para \(p\in\Z\) y \(x>0\),
\begin{equation}
 L_{-p}(x^{-1})=-L_p(x).
 \label{espiral:eq:lp-inversion}
\end{equation}
\end{proposition}
\begin{proof}
Si \(p\neq0\),
\[
 L_{-p}(x^{-1})
 =\frac{(x^{-1})^{-p}-1}{-p}
 =-\frac{x^p-1}{p}.
\]
Para \(p=0\), \(\log(x^{-1})=-\log x\).
\end{proof}

La involución intercambia simultáneamente expansión/contracción y
\(p/-p\). No elimina la memoria de orientación, que permanece en
\(a\).

\section{Límite logarítmico}

Para \(x>0\),
\[
 \lim_{p\to0}\frac{x^p-1}{p}=\log x.
\]
En efecto, \(x^p=\e^{p\log x}=1+p\log x+O(p^2)\). Por tanto la rama
logarítmica no es una excepción añadida: es el límite continuo de la familia.
En la teoría genealógica, sin embargo, \(p=0\) continúa conservando
\(a\) y el ledger; el límite funcional no identifica sus tres
orígenes locales de \cref{espiral:tab:doble-trit}.

\chapter{Clasificación potencia--logarítmica de las curvas publicadas}

\section{Traslación uniforme en la coordenada radial}

Supongamos, después de la publicación del continuo, que
\[
 L_p(r/r_0)=\beta m,
 \qquad
 \theta=\theta_0+\omega m,
 \qquad \omega\neq0.
\]
Eliminando \(m\),
\begin{equation}
 r_p(\theta)=r_0
 \left[1+p\frac{\beta}{\omega}
 (\theta-\theta_0)\right]^{1/p},
 \qquad p\neq0,
 \label{espiral:eq:familia-potencia}
\end{equation}
y
\begin{equation}
 r_0(\theta)=r_0
 \exp\!\left[\frac{\beta}{\omega}
 (\theta-\theta_0)\right].
 \label{espiral:eq:familia-log}
\end{equation}

\begin{theorem}[Publicación de las cinco clases canónicas]
Tras reoriginar \(\theta\) y reescalar \(r\), las especializaciones de
\cref{espiral:eq:familia-potencia,espiral:eq:familia-log} son:
\[
\begin{array}{c|c|c}
p&\text{ecuación normalizada}&\text{reconocimiento convencional}\\ \hline
2&r^2=a^2\theta&\text{espiral de Fermat}\\
1&r=a\theta&\text{espiral de Arquímedes}\\
0&r=r_0\e^{b\theta}&\text{espiral logarítmica}\\
-1&r=a/\theta&\text{espiral hiperbólica}\\
-2&r^2=a^2/\theta&\text{lituus}
\end{array}
\]
en los dominios reales positivos correspondientes.
\end{theorem}
\begin{proof}
Para \(p=1\), \(r\) es afín en \(\theta\), y la reoriginación elimina el
término constante. Para \(p=2\), \(r^2\) es afín. Para \(p=-1\),
\(r^{-1}\) es afín; para \(p=-2\), \(r^{-2}\) es afín. La rama \(p=0\)
es \cref{espiral:eq:familia-log}. Las constantes positivas se absorben en \(a,b\)
sin alterar el exponente.
\end{proof}

\begin{figure}[!htb]
\centering
\begin{tikzpicture}[scale=.62]
\begin{scope}[xshift=0cm,yshift=0cm]
  \draw[->,HMTLine] (-1.9,0)--(1.9,0);
  \draw[->,HMTLine] (0,-1.9)--(0,1.9);
  \draw[HMTBlue,very thick,domain=.15:18.2,samples=180,smooth,variable=\t]
    plot ({sqrt(.14*(\t+1))*cos(deg(\t))},{sqrt(.14*(\t+1))*sin(deg(\t))});
  \node[below,text=HMTGray,font=\scriptsize\sffamily] at (0,-2.1) {Fermat, \(p=2\)};
\end{scope}
\begin{scope}[xshift=5cm,yshift=0cm]
  \draw[->,HMTLine] (-1.9,0)--(1.9,0);
  \draw[->,HMTLine] (0,-1.9)--(0,1.9);
  \draw[HMTNavy,very thick,domain=.15:18.2,samples=180,smooth,variable=\t]
    plot ({.08*(\t+1)*cos(deg(\t))},{.08*(\t+1)*sin(deg(\t))});
  \node[below,text=HMTGray,font=\scriptsize\sffamily] at (0,-2.1) {Arquímedes, \(p=1\)};
\end{scope}
\begin{scope}[xshift=10cm,yshift=0cm]
  \draw[->,HMTLine] (-1.9,0)--(1.9,0);
  \draw[->,HMTLine] (0,-1.9)--(0,1.9);
  \draw[HMTRed,very thick,domain=.15:18.2,samples=180,smooth,variable=\t]
    plot ({.18*exp(.055*\t)*cos(deg(\t))},{.18*exp(.055*\t)*sin(deg(\t))});
  \node[below,text=HMTGray,font=\scriptsize\sffamily] at (0,-2.1) {Logarítmica, \(p=0\)};
\end{scope}
\begin{scope}[xshift=2.5cm,yshift=-5.2cm]
  \draw[->,HMTLine] (-1.9,0)--(1.9,0);
  \draw[->,HMTLine] (0,-1.9)--(0,1.9);
  \draw[HMTGold,very thick,domain=.15:18.2,samples=180,smooth,variable=\t]
    plot ({1.8/(\t+1.8)*cos(deg(\t))},{1.8/(\t+1.8)*sin(deg(\t))});
  \node[below,text=HMTGray,font=\scriptsize\sffamily] at (0,-2.1) {Hiperbólica, \(p=-1\)};
\end{scope}
\begin{scope}[xshift=7.5cm,yshift=-5.2cm]
  \draw[->,HMTLine] (-1.9,0)--(1.9,0);
  \draw[->,HMTLine] (0,-1.9)--(0,1.9);
  \draw[HMTGray,very thick,domain=.15:18.2,samples=180,smooth,variable=\t]
    plot ({sqrt(2.0/(\t+1.8))*cos(deg(\t))},{sqrt(2.0/(\t+1.8))*sin(deg(\t))});
  \node[below,text=HMTGray,font=\scriptsize\sffamily] at (0,-2.1) {Lituus, \(p=-2\)};
\end{scope}
\end{tikzpicture}
\caption{Cinco publicaciones de la familia potencia--logarítmica sobre ejes y escalas comunes.}
\label{espiral:fig:cinco-espirales}
\end{figure}

\section{No exhaustividad y alcance exacto}

El teorema clasifica la familia generada por traslación uniforme en la
coordenada \(L_p\). No clasifica todas las curvas planas posibles. La
taxonomía genealógica completa debe incluir al menos
\[
 (p,a,[\Lambda,C],[\kappa_9],
 \omega_{\mathrm{ang}},\mu_{\mem},\varepsilon_{\mathrm{tw}}).
\]
Dos rutas con el mismo \(p\) pueden diferir en autoescala, orientación,
carry o frontera y publicar curvas distintas. Dos rutas con todos los datos
visibles iguales pueden diferir aún en una coordenada de ledger eliminada.

\section{Holonomía afín general}

Si la coordenada \(u=L_p(r/r_0)\) satisface
\[
 u_{n+1}=\Lambda u_n+C,
\]
la iteración interna exacta, para \(n\in\N\), es
\[
 u_n=\Lambda^n u_0+\sum_{k=0}^{n-1}\Lambda^kC.
\]
Si \(I-\Lambda\) es invertible, puede escribirse
\[
 u_n=u_*+\Lambda^n(u_0-u_*),
 \qquad u_*=(I-\Lambda)^{-1}C.
\]
La interpolación continua sólo se efectúa después de fijar un carácter
escalar positivo \(\chi:U_{\APP}\to\R\) tal que
\(\chi\Lambda=\lambda\chi\), con \(\lambda>0\). Escribiendo
\(u=\chi(u)\), \(C=\chi(C)\), si \(\lambda\neq1\) se tiene
\(u_*=C/(1-\lambda)\). Para \(\theta_n=\theta_0+n\omega\),
\(\omega\neq0\), la realización escalar admite
\begin{equation}
 r(\theta)=r_0\exp_p\!\left[
 u_*+\lambda^{(\theta-\theta_0)/\omega}(u_0-u_*)
 \right].
 \label{espiral:eq:familia-afin-general}
\end{equation}
Las potencias reales pertenecen exclusivamente a esta realización; la
dinámica interna utiliza potencias enteras de \(\Lambda\). La pareja
\((p,[\Lambda,C])\) amplía la lista clásica: \(p\) fija la carta radial y
la clase de conjugación afín fija la dinámica dentro de ella.

\begin{remark}[Punto fijo y electrón]
La existencia de un punto fijo afín no identifica por sí sola ese punto con el
electrón central \((5,5)\). El electrón es un resultado material previo; una
identificación con un punto fijo radial exige un cuadrado de realización
separado.
\end{remark}

\chapter{Lector radial enriquecido}

\section{Obstrucción al lector reducido}

El mapa inicialmente sugerido
\[
 X_{\TPK}^{\enr}\to(p,a,\Lambda,C,m,r,\theta)
\]
mezcla niveles causales y pierde prefijos. Las coordenadas \(r,\theta\) son
arquimedianas, mientras \((\Lambda,C)\) conserva sólo el producto final.

\begin{proposition}[El par final no admite truncamiento natural]
No existe en general una aplicación que recupere el primer factor de una
palabra afín a partir de su producto final.
\end{proposition}
\begin{proof}
En notación \((\Lambda,C)\),
\[
 (1,2)(2,0)=(2,2),
 \qquad
 (1,1)(2,1)=(2,2).
\]
Los dos productos finales coinciden, pero sus primeros prefijos son
\((2,0)\) y \((2,1)\). Una función del producto final tendría que asignarle
dos truncamientos distintos, lo cual es imposible.
\end{proof}

\section{Definición del lector completo}

Sea
\(\widehat\gamma_N=\widehat e_N\cdots\widehat e_1\) una ruta de la
categoría levantada y sea \((H_0,\Led_0)=((I,0),\Led(x_0))\) su dato
inicial. Definimos
\begin{equation}
 \boxed{
 \mathcal P_{\rad,N}^{\enr}(\widehat\gamma_N)
 =\left(
 (H_0,\Led_0);
 \bigl(p(\widehat e_j),a(\widehat e_j),P_j,A_j,
 h_{e_j},H_j,\Led_j(\widehat\gamma_N)\bigr)_{j=1}^{N}
 \right).}
 \label{espiral:eq:lector-radial-enriquecido}
\end{equation}
El codominio \(\mathcal D_{\rad,N}^{\enr}\) contiene secuencias compatibles
de esos datos, incluidas la arista subyacente almacenada en cada entrada del
ledger y sus mapas de fuente y destino. La salida conserva el dato inicial,
cada factor, cada holonomía prefija, los cociclos locales y acumulados y cada
actualización de memoria.

\begin{definition}[Sector radial homogéneo admisible]
Sea \(\mathcal D_{\rad,\infty}^{\mathrm{hom}}\) el espacio de datos marcados
\[
 D=(D^{\enr};r_0,p,u_0,\beta,\theta_0,\omega),
\]
donde \(D^{\enr}\in\mathcal D_{\rad,\infty}^{\enr}\), \(r_0>0\),
\(p\in\{-2,-1,0,1,2\}\), \(u_0,\beta,\theta_0,\omega\in\R\) y
\((\beta,\omega)\neq(0,0)\), tales que
\[
 p(\widehat e_j)=p,\qquad
 u(m)=u_0+\beta m,\qquad
 \theta(m)=\theta_0+\omega m.
\]
La marca forma parte del dato y no se elige después; por ello la evaluación
es una función bien definida.
Su intervalo real de publicación es
\[
 I_D=
 \begin{cases}
 \{m\in\R:1+p(u_0+\beta m)>0\},&p\neq0,\\
 \R,&p=0.
 \end{cases}
\]
Las historias con grado variable permanecen en el ledger enriquecido y se
publican mediante un atlas de cartas por tramos; no determinan una única
curva de grado constante.
\end{definition}

\begin{definition}[Publicación espiral homogénea]
Después de construir el continuo y elegir los caracteres radial y angular
compatibles, se define
\[
 \ev_{\arq}^{\mathrm{esp}}:
 \mathcal D_{\rad,\infty}^{\mathrm{hom}}\longrightarrow
 \coprod_{D}C^\infty(I_D,\R^2)
\]
por
\[
 D\longmapsto
 \left[m\longmapsto
 \bigl(r_0\exp_{p}(u_0+\beta m)\cos(\theta_0+\omega m),
 r_0\exp_{p}(u_0+\beta m)\sin(\theta_0+\omega m)\bigr)\right].
\]
\end{definition}

El lector interno y la evaluación externa están separados. La primera
flecha conserva genealogía; la segunda puede olvidar parte de ella. La
condición \((\beta,\omega)\neq(0,0)\) elimina únicamente la parametrización
constante degenerada; en \(I_D\) la curva publicada es regular.

\section{Conservación como trazabilidad}

La proyección sobre carry, frontera y memoria satisface
\[
 \operatorname{pr}_{\mathrm{carry},frontier,memory}
 \circ\mathcal P_{\rad,N}^{\enr}
 =\Led_{{\mathrm{carry},frontier,memory},N}.
\]
Conservar no significa hacer constantes esas coordenadas. Significa que:
\[
 g(K+9)=g(K),
\]
\[
 q_{\mem}(\Gamma_9x)=q_{\mem}(x)+1,
\]
\[
 B_{K+9}=\operatorname{Tra}_{\Gamma_9}(B_K),
\]
y el carry se actualiza por \cref{espiral:eq:cociclo-carry}. Ninguna actualización
queda oculta por la curva final.

\chapter{Naturalidad, equivariancia y límite inverso}

\section{Naturalidad bajo truncamiento}

Para \(M\le N\), sea
\[
 \pi_{N,M}:\Hist_N^{\enr}(\TPK)\to\Hist_M^{\enr}(\TPK)
\]
el truncamiento de historias, y sea
\[
 q_{N,M}:\mathcal D_{\rad,N}^{\enr}\to
 \mathcal D_{\rad,M}^{\enr}
\]
la eliminación de los datos de índice mayor que \(M\).

\begin{theorem}[Naturalidad del lector radial]
El diagrama
\[
\begin{array}{ccc}
\Hist_N^{\enr}(\TPK)
 & \xrightarrow{\ \mathcal P_{\rad,N}^{\enr}\ }
 & \mathcal D_{\rad,N}^{\enr}\\[1.0em]
{\scriptstyle \pi_{N,M}}\big\downarrow
 && \big\downarrow{\scriptstyle q_{N,M}}\\[1.0em]
\Hist_M^{\enr}(\TPK)
 & \xrightarrow{\ \mathcal P_{\rad,M}^{\enr}\ }
 & \mathcal D_{\rad,M}^{\enr}
\end{array}
\]
conmuta; es decir,
\begin{equation}
 q_{N,M}\circ\mathcal P_{\rad,N}^{\enr}
 =\mathcal P_{\rad,M}^{\enr}\circ\pi_{N,M}.
 \label{espiral:eq:naturalidad-truncamiento}
\end{equation}
\end{theorem}
\begin{proof}
Ambos miembros conservan exactamente el registro inicial y las primeras
\(M\) flechas, junto con sus cociclos locales y acumulados, factores afines,
holonomías prefijas y entradas de ledger de \(\widehat\gamma_N\). No
interviene una reconstrucción desde el producto final: la secuencia completa
forma parte de la definición.
\end{proof}

\section{Lector del límite inverso}

La compatibilidad \cref{espiral:eq:naturalidad-truncamiento} induce
\begin{equation}
 \mathcal P_{\rad,\infty}^{\enr}:
 \varprojlim_N\Hist_N^{\enr}(\TPK)
 \longrightarrow
 \varprojlim_N\mathcal D_{\rad,N}^{\enr}.
 \label{espiral:eq:lector-limite-inverso}
\end{equation}
Cada componente finita se obtiene por truncamiento de la historia completa.
Por ello el lector no inventa una continuidad externa: prolonga de manera
compatible los certificados finitos.

\section{Subdivisión cuádruple}

En el refinamiento \(\sd_4\) del TPK, tres microflechas son quietas y la
cuarta reproduce la transición observable. Si
\[
 \rho_{\rad}(e_1)=\rho_{\rad}(e_2)=\rho_{\rad}(e_3)=I,
 \qquad \rho_{\rad}(e_4)=h_e,
\]
entonces
\[
 H_{\sd_4(e)}=I\,I\,I\,h_e=h_e.
\]
La naturalidad está así certificada para este refinamiento. Para una
sustitución general \(e\mapsto\widetilde e_k\cdots\widetilde e_1\), la
condición exacta es
\[
 h_e=h_{\widetilde e_k}\cdots h_{\widetilde e_1}.
\]
Esta igualdad, y no la semejanza de las curvas publicadas, decide la
compatibilidad de un nuevo refinamiento.

\section{Equivarianza con la holonomía nonádica}

Sea \(b_9(x)\) el bloque de nueve transiciones que realiza \(\Gamma_9\) desde
\(x\), y
\[
 H_9(x)=H_{b_9(x)}=(\Lambda_9(x),C_9(x)).
\]
La acción inducida en el ledger radial añade ese bloque.

\begin{theorem}[Equivarianza de cociclo]
Se cumple
\begin{equation}
 \mathcal P_{\rad}^{\enr}\circ\Gamma_9
 =\widehat\Gamma_9^{\rad}\circ
 \mathcal P_{\rad}^{\enr}.
 \label{espiral:eq:equivarianza-gamma9}
\end{equation}
En componentes,
\[
 P\mapsto P+p(b_9),
 \qquad
 A\mapsto A+a(b_9),
\]
\[
 \Lambda\mapsto\Lambda_9(x)\Lambda,
 \qquad
 C\mapsto C_9(x)+\Lambda_9(x)C,
\]
\[
 g\mapsto g,\qquad q_{\mem}\mapsto q_{\mem}+1.
\]
\end{theorem}
\begin{proof}
La acción de \(\Gamma_9\) concatena \(b_9(x)\) a la historia, donde
\(p(b_9)\) y \(a(b_9)\) son las sumas de los caracteres locales de sus nueve
flechas levantadas. La ley de cociclo suma los grados y la ley afín
\cref{espiral:eq:producto-affend} compone las holonomías. El TPK aporta el retorno de
fase y el incremento de memoria. La definición de
\(\widehat\Gamma_9^{\rad}\) reproduce exactamente esas actualizaciones.
\end{proof}

No es una conmutación literal de endomorfismos del mismo espacio. Los
codominios son distintos y la memoria cambia. La palabra correcta es
\emph{equivarianza}.

\section{Inversión en el sector reversible}

Si \(h_{e^\dagger}=h_e^{-1}\), el diagrama de inversión satisface
\[
 \mathcal P_{\rad}^{\enr}(C_{\mathrm{rev}}\gamma)
 =\mathcal J_{\rad}
 \mathcal P_{\rad}^{\enr}(\gamma),
\]
donde \(\mathcal J_{\rad}\) invierte el orden de factores, reemplaza cada
factor por su inverso y transporta el ledger mediante la involución TPK. En
la salida real, \cref{espiral:eq:lp-inversion} intercambia las dos orientaciones.

\chapter{Clasificación genealógica y no fidelidad de la publicación}

\section{Firma radial plena de una historia}

Para una historia finita
\(\widehat\gamma_N=\widehat e_N\cdots\widehat e_1\), definimos
\begin{equation}
 \mathfrak I_N(\widehat\gamma_N)=
 \left(
 (x_0,H_0,\Led_0);
 \bigl(\widehat e_j,p(\widehat e_j),a(\widehat e_j),P_j,A_j,
 h_{e_j},H_j,\Led_j\bigr)_{j=1}^{N}
 \right).
 \label{espiral:eq:invariante-espiral}
\end{equation}
La firma es plena respecto del lector radial enriquecido escogido: conserva
el objeto inicial, la secuencia tipada de flechas, los caracteres locales y
acumulados, todos los prefijos afines y el ledger completo. No se presenta
como invariante completo de toda la HMT ni sustituye las coordenadas del
estado enriquecido que el lector no tenga por dominio.

\begin{definition}[Equivalencia espiral genealógica]
Dos historias son equivalentes si existe un isomorfismo tipado de sus datos
radiales que transporta término a término la secuencia de
\cref{espiral:eq:invariante-espiral}, respeta fuentes, destinos, composición y
ledgers, y conmuta con todos los truncamientos declarados.
\end{definition}

En el límite inverso se obtiene la familia compatible
\((\mathfrak I_N)_{N\ge0}\). La igualdad de curvas en \(\R^2\) es una
relación más gruesa: la evaluación arquimediana factoriza por la equivalencia
genealógica, pero en general no la refleja.

\section{Clasificación por 3 incrementos visibles}

Sean \(\omega\) el incremento angular, \(\beta\) el incremento radial en
la coordenada \(L_p\) y \(\mu\) el incremento de memoria levantada. La tabla
visible mínima es:
\begin{longtable}{ccccp{.34\textwidth}}
\caption{Clasificación cinemática de las publicaciones.}
\label{espiral:tab:clasificacion-cinematica}\\
\toprule
\(\omega\)&\(\beta\)&\(\mu\)&Tipo publicado&Interpretación\\
\midrule
\endfirsthead
\toprule
\(\omega\)&\(\beta\)&\(\mu\)&Tipo publicado&Interpretación\\
\midrule
\endhead
0&\(\neq0\)&0&recta radial&variación de escala sin giro ni memoria\\
\(\neq0\)&0&0&circunferencia&cierre angular intrínsecamente radial\\
\(\neq0\)&0&\(\neq0\)&hélice cilíndrica&fase periódica con memoria levantada\\
\(\neq0\)&\(\neq0\)&0&espiral plana&variación radial sin eje de memoria visible\\
\(\neq0\)&\(\neq0\)&\(\neq0\)&suspensión helicoidal radial&giro, autoescala y memoria\\
0&0&\(\neq0\)&fibra axial&avance de memoria sin figura transversal\\
\bottomrule
\end{longtable}

Esta tabla clasifica la cinemática publicada, no toda la genealogía. Dos
filas con los mismos tres incrementos pueden diferir en \(p\), carry,
frontera o clase afín.

\section{3 fibras de una misma evaluación visible}

\begin{example}[Circunferencia con memoria nula]
Sea \(H_j=I\), \(\mu=0\) y \(\theta_j=2\pi j/9\). La salida es una
circunferencia y el ledger no registra elevación axial.
\end{example}

\begin{example}[La misma circunferencia como sombra]
Sea \(H_j=I\), \(\mu=1\) y la misma fase angular. En \(\R^3\) se obtiene la
hélice \cref{espiral:eq:inmersion-helice-nonadica}; al aplicar \(\pi_{xy}\), la
salida plana coincide con la del ejemplo anterior.
\end{example}

\begin{example}[Misma espiral, distinto carry]
Sean dos historias con los mismos caracteres reales \(p,\beta,\omega\),
pero con descomposiciones de carry distintas que la evaluación radial envía
al mismo incremento \(\beta\). Ambas publican \(r_p(\theta)\), mientras sus
ledgers difieren. El lector enriquecido las separa; la curva no.
\end{example}

\section{Teorema de clasificación genealógica}

\begin{theorem}[Correspondencias espirales genealógicas]
Fijados el módulo radial universal, la representación local
\(e\mapsto h_e\) y los cociclos radiales del levantamiento TRIT, toda
historia compatible TPK determina:
\begin{enumerate}
\item una secuencia natural de factores y holonomías afines prefijas;
\item una clase genealógica estable bajo truncamiento;
\item una acción equivariante de la holonomía nonádica;
\item en el sector reversible, una historia conjugada de holonomía inversa;
\item para cada sector radial homogéneo admisible, una curva regular en
\(\R^2\) y, al conservar la profundidad, una suspensión en \(\R^3\).
\end{enumerate}
Una historia de grado variable determina en cambio un atlas ordenado de
cartas por tramos. La aplicación a curvas no es fiel: puede identificar
historias con ledgers distintos.
\end{theorem}
\begin{proof}
Los puntos 1 y 2 son los teoremas de holonomía afín y naturalidad. El punto
3 es \cref{espiral:eq:equivarianza-gamma9}; el punto 4 es la inversión en el
subgrupoide reversible. El punto 5 se obtiene restringiendo el límite
inverso al dominio \(\mathcal D_{\rad,\infty}^{\mathrm{hom}}\) y aplicando sus
caracteres radial y angular. La no fidelidad sigue de la
\cref{espiral:fig:circulo-helice} y de la proposición general del primer capítulo.
\end{proof}

\begin{statusbox}[title={Hipótesis, dominio y alcance del teorema}]
La holonomía nonádica, la hélice, el tornillo de Pell y los cociclos del TPK
son resultados recuperados. La envolvente afín, el lector por prefijos, su
naturalidad y la publicación \(L_p\) constituyen una formalización nueva
exacta en los dominios declarados. La carta local
\(p(\widehat e)=\tau_+(e)+\tau_-(e)\) y
\(a(\widehat e)=\tau_+(e)-\tau_-(e)\) es exacta en la categoría levantada
de 2 semicoronas; su empleo fuera de ese dominio exige exhibir el funtor que
aporte ambas orientaciones. La publicación como una sola curva se restringe
al sector homogéneo admisible definido en este capítulo.
\end{statusbox}
